% ============================================================
% CUSTOM COMMANDS — Lệnh tắt & tiện ích
% ============================================================

% --- Counters ---
\newcounter{baitoanCounter}[chapter]
\renewcommand{\thebaitoanCounter}{\thechapter.\arabic{baitoanCounter}}

% --- Đáp án ABCD ---
\newcommand{\dapan}[4]{%
  \begin{enumerate}[label=\textbf{\Alph*.}, itemsep=2pt, leftmargin=2em]
    \item #1
    \item #2
    \item #3
    \item #4
  \end{enumerate}
}

% Đáp án 2 cột
\newcommand{\dapanhai}[4]{%
  \begin{multicols}{2}
    \begin{enumerate}[label=\textbf{\Alph*.}, itemsep=2pt, leftmargin=2em]
      \item #1
      \item #2
      \item #3
      \item #4
    \end{enumerate}
  \end{multicols}
}

% Đáp án 4 cột
\newcommand{\dapanbon}[4]{%
  \begin{multicols}{4}
    \begin{enumerate}[label=\textbf{\Alph*.}, itemsep=2pt, leftmargin=2em]
      \item #1
      \item #2
      \item #3
      \item #4
    \end{enumerate}
  \end{multicols}
}

% --- Đánh dấu đáp án đúng ---
\newcommand{\dung}[1]{\textbf{\textcolor{ForumGreen}{#1}}}



% --- Separator ---
\newcommand{\separator}{%
  \begin{center}
    \textcolor{ForumGray}{\rule{0.6\textwidth}{0.5pt}}
  \end{center}
}

% --- Toán học tiện ích ---
\newcommand{\R}{\mathbb{R}}
\newcommand{\N}{\mathbb{N}}
\newcommand{\Z}{\mathbb{Z}}
\newcommand{\Q}{\mathbb{Q}}
\newcommand{\C}{\mathbb{C}}
\DeclareMathOperator{\Ima}{Im}
\DeclareMathOperator{\Rea}{Re}
